\documentclass[../../main.tex]{subfiles}
\begin{document}
\begin{oframed}
  \subfile{problem}
\end{oframed}

{\bf [解答]}

\begin{figure}[htb]
  \centering
  \begin{tikzpicture}[scale=1.5, >=stealth, every node/.style={font=\small}]
    \coordinate (O) at (0, 0);

    % 固定正方形 B: [0, 2] x [1, 3], 中心 (1, 2)
    \coordinate (B_BL) at (0, 1);
    \coordinate (B_BR) at (2, 1);
    \coordinate (B_TR) at (2, 3);
    \coordinate (B_TL) at (0, 3);
    \coordinate (C_B) at (1, 2);

    % 動く正方形 A (最大値 theta = pi/3 の位置)
    % 中心 P(cos 60°, sin 60°) = (0.5, 0.866)
    \def\cVal{0.5}
    \def\sVal{0.8660}
    \coordinate (P) at (\cVal, \sVal);
    \coordinate (A_BL) at ({\cVal - 1}, {\sVal - 1});
    \coordinate (A_BR) at ({\cVal + 1}, {\sVal - 1});
    \coordinate (A_TR) at ({\cVal + 1}, {\sVal + 1});
    \coordinate (A_TL) at ({\cVal - 1}, {\sVal + 1});

    % 共通部分の斜線ハッチング: [0, cos theta + 1] x [1, sin theta + 1]
    \fill[pattern=north east lines, pattern color=blue!50]
    (0, 1) -- ({\cVal + 1}, 1) -- ({\cVal + 1}, {\sVal + 1}) -- (0, {\sVal + 1}) -- cycle;

    % 座標軸
    \draw[->, thick, gray!80!black] (-1.8, 0) -- (2.6, 0) node[right] {$x$};
    \draw[->, thick, gray!80!black] (0, -1.4) -- (0, 3.4) node[above] {$y$};
    \node[below left=2pt] at (O) {$O$};

    % 単位円 x^2 + y^2 = 1
    \draw[line width=1.0pt, gray!60!black, dashed] (O) circle (1);
    \node[below left, font=\footnotesize, gray!70!black] at (-1, 0) {$-1$};
    \node[below, font=\footnotesize, gray!70!black] at (1, 0) {$1$};

    % 正方形 B の輪郭
    \draw[line width=1.1pt, gray!85!black] (B_BL) -- (B_BR) -- (B_TR) -- (B_TL) -- cycle;
    \node[above right, font=\footnotesize, gray!85!black] at (B_TR) {正方形 $B$};
    \node[circle, fill=gray!80!black, inner sep=1.2pt] at (C_B) {};
    \node[right=2pt, font=\footnotesize] at (C_B) {$(1, 2)$};

    % 正方形 A の輪郭
    \draw[line width=1.2pt, blue!85!black] (A_BL) -- (A_BR) -- (A_TR) -- (A_TL) -- cycle;
    \node[above left, font=\footnotesize, blue!85!black] at (A_TL) {正方形 $A$};

    % 共通部分の輪郭線
    \draw[line width=1.3pt, red!85!black]
    (0, 1) -- ({\cVal + 1}, 1) -- ({\cVal + 1}, {\sVal + 1}) -- (0, {\sVal + 1}) -- cycle;
    \node[red!85!black, font=\footnotesize] at (0.75, 1.45) {$S(\theta)$};

    % 中心 P と動径 OP
    \draw[line width=1.0pt, blue!85!black] (O) -- (P);
    \draw[->, gray!70!black] (0.35, 0) arc (0:60:0.35);
    \node[font=\footnotesize, gray!80!black] at (25:0.5) {$\theta$};
    \node[circle, fill=blue!85!black, inner sep=1.3pt] at (P) {};
    \node[below right=-1pt, font=\footnotesize, blue!85!black] at (P) {$P(\cos\theta,\sin\theta)$};

    % 軸への点線投影
    \draw[densely dotted, thick, gray!70!black] ({\cVal + 1}, 0) -- ({\cVal + 1}, {\sVal + 1});
    \node[below=2pt, font=\footnotesize] at ({\cVal + 1}, 0) {$\cos\theta+1$};

    \draw[densely dotted, thick, gray!70!black] (0, 1) -- (B_BL);
    \node[left=2pt, font=\footnotesize] at (0, 1) {$1$};

    \draw[densely dotted, thick, gray!70!black] (0, {\sVal + 1}) -- ({\cVal + 1}, {\sVal + 1});
    \node[left=2pt, font=\footnotesize] at (0, {\sVal + 1}) {$\sin\theta+1$};
  \end{tikzpicture}
  \caption{正方形 $A$ と固定正方形 $B$ の共通部分 $S(\theta)$}
  \label{fig:squares_intersection}
\end{figure}

$A$の中心$(c,s) \ (0\le\theta<2\pi)$とすると，$A$の周及び内部の点は不等式
\begin{align}
  \begin{cases}
    \cos\theta-1\le x\le \cos\theta+1 \\
    \sin\theta-1\le y\le \sin\theta+1
  \end{cases} \label{eq:1}
\end{align}
となる．一方$B$の周及び内部は不等式
\begin{align}
  \begin{cases}
    0\le x\le2 \\
    1\le y\le3
  \end{cases}\label{eq:2}
\end{align}
で表される．\cref{eq:1,eq:2}が共通部分を持つための条件は$\sin\theta\ge0$つまり$0\le\theta\le\pi$の時で，このもとで共通部分は
\begin{align}
  \begin{cases}
    0\le x\le \cos\theta+1\\
    1\le y\le \sin\theta+1
  \end{cases}\label{eq:3}
\end{align}
で表される長方形である．故にこの面積$S(\theta)$は
\begin{align}
  S(\theta)=\sin\theta(\cos\theta+1)
\end{align}
である．このとき $S'(\theta)=(2\cos\theta-1)(\cos\theta+1)$ であるから，\cref{tab:s_theta_increments}の増減表を得る．

\begin{table}[htb]
  \centering
  \caption{$S(\theta)$ の増減表}
  \label{tab:s_theta_increments}
  \begin{tabular}{|c|c|c|c|c|c|}
    \hline
    $\theta$ & $0$ & $\cdots$ & $\dfrac{\pi}{3}$ & $\cdots$ & $\pi$ \\ \hline
    $\cos\theta$ & $1$ & $\cdots$ & $\dfrac{1}{2}$ & $\cdots$ & $-1$ \\ \hline
    $S'(\theta)$ & & $+$ & $0$ & $-$ & $0$ \\ \hline
    $S(\theta)$ & $0$ & $\nearrow$ & $\dfrac{3\sqrt{3}}{4}$ & $\searrow$ & $0$ \\ \hline
  \end{tabular}
\end{table}
従って$\max S(\theta)=S(\pi/3)=\dfrac{3\sqrt{3}}{4}$である．$\cdots$(答)
\newpage
\end{document}
